\chapter{I Numeri Reali e il Calcolo Letterale}

\section{I Numeri Reali}
\subsection{Definizione e struttura di $\mathbb{R}$}
\subsubsection{Numeri razionali e irrazionali}
L'insieme dei numeri razionali $\mathbb{Q}$ è costituito dalle classi di equivalenza delle frazioni $\frac{p}{q}$ con $p,q \in \mathbb{Z}$ e $q \neq 0$. Sebbene $\mathbb{Q}$ sia un campo totalmente ordinato e denso, esso è \textit{incompleto}. L'equazione $x^2 = 2$, ad esempio, non ammette soluzioni in $\mathbb{Q}$.
I numeri non esprimibili come rapporto di interi, che possiedono un allineamento decimale illimitato e non periodico, sono detti \textit{irrazionali} (es. $\sqrt{2}, \pi, e$). L'unione degli insiemi dei numeri razionali e irrazionali costituisce l'insieme dei numeri reali $\mathbb{R}$.

\textbf{Esempio.} Si dimostra per assurdo che $\sqrt{2} \notin \mathbb{Q}$. Se per assurdo $\sqrt{2} = \frac{p}{q}$ con $\text{MCD}(p,q)=1$, elevando al quadrato si ottiene $2q^2 = p^2$. Dunque $p^2$ è pari, il che implica che anche $p$ sia pari (ovvero $p=2k$). Sostituendo si ha $2q^2 = 4k^2 \implies q^2 = 2k^2$. Ne consegue che anche $q$ è pari, contraddicendo l'ipotesi iniziale di coprimalità tra $p$ e $q$.

\subsubsection{Assioma di continuità (Completezza di $\mathbb{R}$)}
Ciò che distingue strutturalmente $\mathbb{R}$ da $\mathbb{Q}$ è l'Assioma di Completezza (o di Dedekind). 
\textbf{Definizione (Assioma di separazione):} Dati due sottoinsiemi non vuoti $A, B \subset \mathbb{R}$ tali che per ogni $a \in A$ e per ogni $b \in B$ risulti $a \le b$ (insiemi separati), esiste sempre almeno un elemento $c \in \mathbb{R}$ tale che $a \le c \le b$ per ogni $a \in A$ e $b \in B$.
L'elemento $c$ è detto \textit{elemento separatore}. In $\mathbb{R}$ esso esiste sempre; in $\mathbb{Q}$ no. 

\textbf{Esempio.} Siano $A=\{x \in \mathbb{Q} \mid x>0 \land x^2<2\}$ e $B=\{x \in \mathbb{Q} \mid x>0 \land x^2>2\}$. $A$ e $B$ sono separati, ma non ammettono alcun elemento separatore in $\mathbb{Q}$. In $\mathbb{R}$, l'elemento separatore è $\sqrt{2}$.

\subsubsection{Intervalli e intorni}
\textbf{Definizione:} Un intervallo è un sottoinsieme di $\mathbb{R}$ che gode della proprietà di connessione: se contiene due numeri, contiene anche tutti i numeri compresi tra essi. 
Tipi di intervalli limitati (con estremi $a, b \in \mathbb{R}$ e $a < b$):
\begin{itemize}
    \item Chiuso: $[a,b] = \{x \in \mathbb{R} \mid a \le x \le b\}$
    \item Aperto: $(a,b) = \{x \in \mathbb{R} \mid a < x < b\}$
\end{itemize}
\textbf{Definizione:} Si definisce \textit{intorno circolare} (o sferico) di un punto $x_0 \in \mathbb{R}$ di raggio $r > 0$ l'intervallo aperto $I_r(x_0) = (x_0 - r, x_0 + r)$. Tale concetto è fondamentale per la successiva formalizzazione topologica dei limiti.

\subsection{Valore assoluto}
\subsubsection{Definizione e proprietà geometriche}
\textbf{Definizione:} Il valore assoluto (o modulo) di un numero reale $x$ è definito come:
\[
|x| = \begin{cases} 
x & \text{se } x \ge 0 \\ 
-x & \text{se } x < 0 
\end{cases}
\]
Geometricamente, rappresenta la distanza dell'origine dal punto di coordinata $x$ sulla retta reale.
\textbf{Proprietà principali:}
\begin{itemize}
    \item $|x| \ge 0$ e $|x| = 0 \iff x = 0$ (Definita positività)
    \item $|xy| = |x||y|$ (Moltiplicatività)
    \item $|x+y| \le |x| + |y|$ (Disuguaglianza triangolare)
\end{itemize}

\subsubsection{Distanza tra due punti sulla retta reale}
\textbf{Definizione:} La distanza $d(x,y)$ tra due punti $x, y \in \mathbb{R}$ è definita come $d(x,y) = |x-y|$. 
Questa funzione soddisfa gli assiomi di metrica (positività, simmetria, disuguaglianza triangolare), dotando $\mathbb{R}$ della struttura di spazio metrico unidimensionale.

\section{Monomi e Polinomi}
\subsection{Monomi}
\subsubsection{Definizione, grado e monomi simili}
\textbf{Definizione:} Un monomio è un'espressione algebrica in cui compaiono esclusivamente operazioni di moltiplicazione tra una parte numerica (coefficiente) e variabili letterali con esponenti interi non negativi. La forma normale è $k \cdot x_1^{a_1} x_2^{a_2} \dots x_n^{a_n}$.
\textbf{Definizione:} Il \textit{grado complessivo} di un monomio è la somma degli esponenti di tutte le sue variabili.
\textbf{Definizione:} Due monomi si dicono \textit{simili} se e solo se presentano la medesima parte letterale.

\subsubsection{Operazioni tra monomi}
\begin{itemize}
    \item \textbf{Addizione/Sottrazione:} Operazione definita internamente solo tra monomi simili. Si sommano algebricamente i coefficienti. \textbf{Esempio:} $3x^2y - 5x^2y = -2x^2y$.
    \item \textbf{Moltiplicazione:} Prodotto dei coefficienti e somma degli esponenti per le variabili omonime. \textbf{Esempio:} $(2x^3y) \cdot (-3xy^2) = -6x^4y^3$.
    \item \textbf{Divisione:} Il quoziente è un monomio se e solo se il grado di ogni variabile nel dividendo è maggiore o uguale a quello corrispondente nel divisore.
    \item \textbf{Potenza:} Si eleva il coefficiente all'esponente richiesto e si moltiplicano gli esponenti delle variabili per l'esponente della potenza.
\end{itemize}

\subsection{Polinomi}
\subsubsection{Definizione e grado di un polinomio}
\textbf{Definizione:} Un polinomio è la somma algebrica di un numero finito di monomi. Se i monomi (termini) non sono simili tra loro, il polinomio si dice ridotto in \textit{forma normale}.
\textbf{Definizione:} Il grado complessivo di un polinomio è il massimo tra i gradi complessivi dei monomi non nulli che lo compongono.

\subsubsection{Somma algebrica e prodotto di polinomi}
L'insieme dei polinomi in una variabile a coefficienti reali, denotato con $\mathbb{R}[x]$, possiede una struttura algebrica di anello commutativo con unità rispetto alle operazioni di addizione e moltiplicazione.
\begin{itemize}
    \item \textbf{Somma:} Consiste nella somma formale e successiva riduzione dei termini simili.
    \item \textbf{Prodotto:} Si applica reiteratamente la proprietà distributiva, moltiplicando ogni termine del primo polinomio per ogni termine del secondo.
\end{itemize}

\subsection{Prodotti notevoli}
I prodotti notevoli sono identità algebriche fondamentali che permettono di calcolare specifici prodotti o potenze in modo diretto.
\subsubsection{Quadrato e cubo di binomio}
\begin{itemize}
    \item $(a+b)^2 = a^2 + 2ab + b^2$
    \item $(a+b)^3 = a^3 + 3a^2b + 3ab^2 + b^3$
\end{itemize}
\subsubsection{Somma per differenza}
\begin{itemize}
    \item $(a+b)(a-b) = a^2 - b^2$
\end{itemize}
\textbf{Esempio:} $(2x - 3y)(2x + 3y) = 4x^2 - 9y^2$.
\subsubsection{Quadrato di trinomio}
\begin{itemize}
    \item $(a+b+c)^2 = a^2 + b^2 + c^2 + 2ab + 2ac + 2bc$
\end{itemize}
\subsubsection{Potenza $n$-esima di un binomio (Triangolo di Tartaglia)}
\textbf{Teorema del Binomio di Newton:} Per ogni $n \in \mathbb{N}$, vale lo sviluppo:
\[
(a+b)^n = \sum_{k=0}^n \binom{n}{k} a^{n-k}b^k
\]
I coefficienti binomiali $\binom{n}{k}$ possono essere ricavati analiticamente tramite i fattoriali oppure disposti ricorsivamente nel Triangolo di Tartaglia.

\subsection{Divisione tra polinomi}
\subsubsection{Divisione euclidea}
\textbf{Teorema:} Dati due polinomi $A(x), B(x) \in \mathbb{R}[x]$ con $B(x) \neq 0$, esistono e sono unici due polinomi $Q(x)$ (quoziente) e $R(x)$ (resto) tali che:
\[ A(x) = B(x)Q(x) + R(x) \]
con il grado di $R(x)$ strettamente minore del grado di $B(x)$, oppure $R(x)=0$ (polinomio nullo).

\subsubsection{Regola di Ruffini e Teorema del Resto}
La Regola di Ruffini è un algoritmo sintetico per la divisione euclidea, applicabile esclusivamente quando il divisore è un binomio monico di primo grado della forma $(x-c)$.
\textbf{Teorema del Resto:} Il resto della divisione di un polinomio $P(x)$ per il binomio $(x-c)$ è dato dal valore che il polinomio assume quando si valuta la variabile $x$ in $c$, ovvero $R = P(c)$.
\textbf{Esempio:} Il resto della divisione di $P(x) = x^3 - 2x + 5$ per $(x-2)$ è $P(2) = 2^3 - 2(2) + 5 = 9$.

\subsubsection{Teorema di Ruffini (zeri di un polinomio)}
\textbf{Teorema (Corollario del Teorema del Resto):} Un polinomio $P(x)$ è divisibile per $(x-c)$ (cioè il resto è nullo) se e solo se $P(c) = 0$. In tal caso, $c$ è detto \textit{radice} o \textit{zero} del polinomio.

\section{Scomposizioni e Frazioni Algebriche}
\subsection{Scomposizione in fattori}
Scomporre un polinomio significa fattorizzarlo, ovvero esprimerlo come prodotto di polinomi di grado inferiore. Un polinomio si dice \textit{irriducibile} in $\mathbb{R}[x]$ se non ammette fattorizzazioni non banali.

\subsubsection{Raccoglimento totale a fattore comune}
Sfrutta in senso inverso la proprietà distributiva: si individua il M.C.D. tra tutti i termini del polinomio e lo si raccoglie a fattore.
\textbf{Esempio:} $6x^3y^2 - 9x^2y^3 = 3x^2y^2(2x - 3y)$.

\subsubsection{Raccoglimento parziale}
Si effettua tipicamente in polinomi con un numero pari di termini, raccogliendo fattori comuni a sottogruppi, per poi procedere con un raccoglimento totale sulla parentesi binomiale ottenuta.
\textbf{Esempio:} $ax + ay + 2x + 2y = a(x+y) + 2(x+y) = (a+2)(x+y)$.

\subsubsection{Scomposizione mediante prodotti notevoli}
Si basa sul riconoscimento visivo degli sviluppi dei prodotti notevoli analizzati in precedenza.
\textbf{Esempio:} $16x^4 - 1 = (4x^2 - 1)(4x^2 + 1) = (2x - 1)(2x + 1)(4x^2 + 1)$. (Notare che $4x^2+1$ è irriducibile in $\mathbb{R}$).

\subsubsection{Scomposizione del trinomio di secondo grado (trinomio speciale)}
Un trinomio monico del tipo $x^2 + sx + p$ si scompone in $(x-x_1)(x-x_2)$ qualora si individuino due numeri $x_1, x_2$ la cui somma sia $s$ e il cui prodotto sia $p$.
Nel caso generale $ax^2 + bx + c$, se l'equazione associata ammette radici $x_1, x_2$ reali (cioè se il discriminante $\Delta \ge 0$), la scomposizione è $a(x-x_1)(x-x_2)$. Se $\Delta < 0$, il trinomio è irriducibile in $\mathbb{R}$.

\subsubsection{Scomposizione mediante la Regola di Ruffini}
Se un polinomio $P(x)$ a coefficienti interi ammette una radice razionale $\frac{p}{q}$, allora $p$ divide il termine noto e $q$ divide il coefficiente direttore (Teorema delle radici razionali). Individuata una radice $c$, $P(x)$ si fattorizza come $(x-c)Q(x)$, abbassando il grado del polinomio tramite l'algoritmo di Ruffini.

\subsection{M.C.D. e m.c.m. di polinomi}
\begin{itemize}
    \item \textbf{M.C.D.:} È il polinomio monico di grado massimo che divide tutti i polinomi dati. Operativamente, si scompongono i polinomi in fattori irriducibili e si moltiplicano i soli fattori \textit{comuni}, presi una singola volta con l'esponente \textit{minore}.
    \item \textbf{m.c.m.:} È il polinomio monico di grado minimo che è multiplo di tutti i polinomi dati. Si moltiplicano i fattori \textit{comuni e non comuni}, presi una singola volta con l'esponente \textit{maggiore}.
\end{itemize}

\subsection{Frazioni algebriche}
\subsubsection{Condizioni di esistenza}
\textbf{Definizione:} Una frazione algebrica è definita come il rapporto $\frac{A(x)}{B(x)}$ tra due polinomi, con $B(x)$ non nullo. Il dominio (o Condizione di Esistenza, C.E.) si determina ponendo il denominatore diverso da zero: $B(x) \neq 0$.
\textbf{Esempio:} Per $\frac{x+1}{x^2 - 4}$, la C.E. è $x^2 - 4 \neq 0 \implies x \neq \pm 2$.

\subsubsection{Semplificazione di frazioni algebriche}
In virtù della proprietà invariantiva, la divisione del numeratore e del denominatore per un loro fattore comune genera una frazione equivalente. La semplificazione richiede rigorosamente la preliminare fattorizzazione completa di $A(x)$ e $B(x)$.

\subsubsection{Operazioni (addizione, moltiplicazione, divisione, potenza)}
\begin{itemize}
    \item \textbf{Addizione algebrica:} Si determinano le C.E., si calcola il m.c.m. dei denominatori scomposti e si riducono le frazioni allo stesso denominatore prima di sommare i numeratori.
    \item \textbf{Moltiplicazione:} È buona prassi eseguire preventivamente semplificazioni "a croce" tra i numeratori e i denominatori fattorizzati, per poi effettuare il prodotto in riga.
    \item \textbf{Divisione:} Si riconduce a una moltiplicazione della prima frazione per la frazione reciproca del divisore.
\end{itemize}


\section{Esercizi}

\subsection{Struttura di $\mathbb{R}$, Assiomi e Valore Assoluto}
\begin{enumerate}
    \item \tagprob{Numeri Reali}{2} Dimostra per assurdo che $\sqrt{3} \notin \mathbb{Q}$.
    \item \tagprob{Numeri Reali}{2} Dimostra per assurdo che $\sqrt{5} \notin \mathbb{Q}$.
    \item \tagprob{Numeri Reali}{3} Sapendo che $\sqrt{2} \notin \mathbb{Q}$, dimostra per assurdo che anche $\sqrt{2} + 1 \notin \mathbb{Q}$.
    \item \tagprob{Numeri Reali}{3} Dimostra che la somma di un numero razionale non nullo e un numero irrazionale è sempre un numero irrazionale.
    \item \tagprob{Assioma Completezza}{2} Dati $A=\{x \in \mathbb{Q} \mid x>0, x^2 < 3\}$ e $B=\{x \in \mathbb{Q} \mid x>0, x^2 > 3\}$, individua l'elemento separatore in $\mathbb{R}$ e dimostra che $A$ e $B$ sono separati.
    \item \tagprob{Assioma Completezza}{2} Dati $A=\{x \in \mathbb{Q} \mid x^3 < 2\}$ e $B=\{x \in \mathbb{Q} \mid x^3 > 2\}$, individua l'elemento separatore in $\mathbb{R}$.
    \item \tagprob{Assioma Completezza}{3} Spiega perché l'insieme $\mathbb{N}$ non ammette l'assioma di completezza (fornisci un controesempio con insiemi separati).
    \item \tagprob{Intorni}{1} Scrivi sotto forma di intervallo aperto l'intorno circolare $I_{0.1}(5)$ e calcolane l'ampiezza.
    \item \tagprob{Intorni}{1} Determina il centro $x_0$ e il raggio $r$ dell'intorno circolare corrispondente all'intervallo aperto $(-2, 6)$.
    \item \tagprob{Intorni}{2} Interseca gli intorni $I_2(0)$ e $I_3(1)$ ed esprimi il risultato come intervallo.
    \item \tagprob{Valore Assoluto}{1} Scrivi l'insieme dei punti la cui distanza da $x_0 = -1$ è pari a $5$.
    \item \tagprob{Valore Assoluto}{1} Scrivi l'insieme dei punti la cui distanza da $x_0 = 3$ è strettamente minore di $2$.
    \item \tagprob{Valore Assoluto}{2} Risolvi usando il concetto di distanza sulla retta reale: $|x - 4| \le 1$.
    \item \tagprob{Valore Assoluto}{3} Sfruttando la disuguaglianza triangolare, dimostra che $|a-b| \le |a-c| + |c-b|$ per ogni $a,b,c \in \mathbb{R}$.
\end{enumerate}

\subsection{Monomi: Operazioni e Semplificazioni}
\begin{multicols}{2}
\begin{enumerate}
    \item \tagprob{Monomi}{1} $(3x^2y)(-2xy^3)$
    \item \tagprob{Monomi}{1} $\left(-\frac{1}{2}a^3b^2\right)\left(\frac{4}{3}ab\right)$
    \item \tagprob{Monomi}{1} $(-5x^4y^2)^3$
    \item \tagprob{Monomi}{1} $\left(\frac{2}{3}m^2n^3\right)^2 \cdot \left(-\frac{9}{4}mn\right)$
    \item \tagprob{Monomi}{1} $(12x^5y^3) \div (-3x^2y)$
    \item \tagprob{Monomi}{2} $\left(-\frac{3}{5}a^4b^2\right)^3 \div \left(\frac{9}{25}a^6b^4\right)$
    \item \tagprob{Monomi}{2} $[(-2x^2y)^3 \cdot (xy^2)^2] \div (-4x^4y^5)$
    \item \tagprob{Monomi}{2} $\left(\frac{1}{8}a^3b^3\right) \div \left(-\frac{1}{2}ab\right)^3$
    \item \tagprob{Monomi}{2} $3x^2y(2xy) - 5x(x^2y^2) + (-2x^3y^2)$
    \item \tagprob{Monomi}{3} Espressione: $\left[ \left(-\frac{1}{2}x^n y^{n+1}\right)^2 \cdot (4x^2y) \right] \div (x^{2n}y^{2n})$ \textit{(con $n \in \mathbb{N}$)}
\end{enumerate}
\end{multicols}

\subsection{Polinomi e Prodotti Notevoli}
\begin{multicols}{2}
\begin{enumerate}
    \item \tagprob{Polinomi}{1} Svolgi: $2x(x^2 - 3x + 1) - x^2(2x - 3)$
    \item \tagprob{Polinomi}{1} Svolgi: $(a+b)(a^2-ab+b^2)$
    \item \tagprob{Prod Notevole}{1} Sviluppa: $(3x - 2y)^2$
    \item \tagprob{Prod Notevole}{1} Sviluppa: $\left(\frac{1}{2}a + 4b\right)^2$
    \item \tagprob{Prod Notevole}{1} Sviluppa: $(-x^2 - 3y)^2$
    \item \tagprob{Prod Notevole}{1} Calcola: $(5a - 2b)(5a + 2b)$
    \item \tagprob{Prod Notevole}{1} Calcola: $(x^2 + \frac{1}{3}y)(x^2 - \frac{1}{3}y)$
    \item \tagprob{Prod Notevole}{2} Sviluppa: $(2x - y + 1)^2$
    \item \tagprob{Prod Notevole}{2} Sviluppa: $(a^2 - 2a + 3)^2$
    \item \tagprob{Prod Notevole}{2} Sviluppa: $(x - 2y)^3$
    \item \tagprob{Prod Notevole}{2} Sviluppa: $\left(2a + \frac{1}{2}\right)^3$
    \item \tagprob{Prod Notevole}{2} Sviluppa: $(-x^2 - y)^3$
    \item \tagprob{Tartaglia}{2} Sviluppa: $(a - b)^4$
    \item \tagprob{Tartaglia}{2} Sviluppa: $(2x + y)^5$
    \item \tagprob{Tartaglia}{3} Determina il coefficiente di $x^4y^2$ nello sviluppo di $(x - 2y)^6$
\end{enumerate}
\end{multicols}

\subsection{Divisione Euclidea, Teorema del Resto e Ruffini}
\begin{multicols}{2}
\begin{enumerate}
    \item \tagprob{Div Euclidea}{1} $(x^3 - 4x^2 + 5x - 2) \div (x^2 - x + 1)$
    \item \tagprob{Div Euclidea}{1} $(2x^4 - 3x^3 + 5x - 1) \div (x^2 + 2)$
    \item \tagprob{Div Euclidea}{1} $(x^5 + 3x^2 - 4) \div (x^3 - 1)$
    \item \tagprob{Div Euclidea}{2} $(3x^4 - x^2 + 2x - 5) \div (x^2 - 2x + 3)$
    \item \tagprob{Div Euclidea}{2} $(a^4 - b^4) \div (a^2 + b^2)$
    \item \tagprob{Teorema Resto}{1} Calcola il resto: $(x^3 - 2x^2 + 4x - 5) \div (x - 2)$
    \item \tagprob{Teorema Resto}{1} Calcola il resto: $(2x^4 + 3x^3 - x + 1) \div (x + 1)$
    \item \tagprob{Teorema Resto}{2} Determina $k$ affinché $P(x) = x^3 - kx^2 + x + 6$ sia divisibile per $(x - 3)$.
    \item \tagprob{Teorema Resto}{2} Determina $k$ affinché il resto di $(2x^3 + kx - 4) \div (x + 2)$ sia pari a $10$.
    \item \tagprob{Ruffini}{1} Dividi con Ruffini: $(x^3 - 6x^2 + 11x - 6) \div (x - 1)$
    \item \tagprob{Ruffini}{1} Dividi con Ruffini: $(2x^4 - 5x^3 + 3x - 2) \div (x - 2)$
    \item \tagprob{Ruffini}{1} Dividi con Ruffini: $(x^5 + 32) \div (x + 2)$
    \item \tagprob{Ruffini}{2} Dividi con Ruffini: $(3x^3 - 4x^2 - 5x + 2) \div \left(x - \frac{1}{3}\right)$
    \item \tagprob{Ruffini}{2} Dividi con Ruffini (esteso): $(2x^3 - x^2 - x + 3) \div (2x + 1)$
    \item \tagprob{Ruffini}{3} Determina per quali $n \in \mathbb{N}$ il polinomio $x^n + 1$ è divisibile per $x+1$.
\end{enumerate}
\end{multicols}

\subsection{Scomposizione in Fattori}
\textit{Fattorizza completamente i seguenti polinomi in $\mathbb{R}[x]$:}
\begin{multicols}{2}
\begin{enumerate}
    % Raccoglimento Totale e Parziale (10 varianti)
    \item \tagprob{Raccoglimento}{1} $12a^3b^2 - 18a^2b^3 + 6a^2b^2$
    \item \tagprob{Raccoglimento}{1} $5x^4y - 10x^3y^2 + 15x^2y^3$
    \item \tagprob{Raccoglimento}{1} $x(a-b) - y(a-b)$
    \item \tagprob{Raccoglimento}{1} $a(x+1) + 2(x+1)^2$
    \item \tagprob{Raccoglimento}{2} $ax - ay + bx - by$
    \item \tagprob{Raccoglimento}{2} $x^3 - 2x^2 + 3x - 6$
    \item \tagprob{Raccoglimento}{2} $2a^3 - 3a^2b + 4ab^2 - 6b^3$
    \item \tagprob{Raccoglimento}{2} $x^2y - x^2 + y - 1$
    \item \tagprob{Raccoglimento}{2} $5ax - 5ay - 2bx + 2by$
    \item \tagprob{Raccoglimento}{3} $x^{n+2} - x^{n+1} + x^n$ \textit{(con $n \ge 1$)}
    
    % Prodotti Notevoli (10 varianti)
    \item \tagprob{Scomp Notevole}{1} $x^2 - 16y^2$
    \item \tagprob{Scomp Notevole}{1} $25a^4 - 9b^2$
    \item \tagprob{Scomp Notevole}{1} $x^2 + 6x + 9$
    \item \tagprob{Scomp Notevole}{1} $4a^2 - 12ab + 9b^2$
    \item \tagprob{Scomp Notevole}{2} $a^3 - 8$
    \item \tagprob{Scomp Notevole}{2} $27x^3 + 1$
    \item \tagprob{Scomp Notevole}{2} $x^4 - 81$
    \item \tagprob{Scomp Notevole}{2} $16x^4 - y^4$
    \item \tagprob{Scomp Notevole}{3} $a^2 - b^2 - 2a + 1$
    \item \tagprob{Scomp Notevole}{3} $x^2 - y^2 + 4x + 4$
    
    % Trinomio Speciale (10 varianti)
    \item \tagprob{Scomp Trinomio}{1} $x^2 - 5x + 6$
    \item \tagprob{Scomp Trinomio}{1} $x^2 + 7x + 10$
    \item \tagprob{Scomp Trinomio}{1} $x^2 - x - 12$
    \item \tagprob{Scomp Trinomio}{1} $x^2 + 3x - 18$
    \item \tagprob{Scomp Trinomio}{1} $x^2 - 8x + 15$
    \item \tagprob{Scomp Trinomio}{2} $2x^2 - 5x + 3$
    \item \tagprob{Scomp Trinomio}{2} $3x^2 + 5x - 2$
    \item \tagprob{Scomp Trinomio}{2} $4x^2 - 4x - 3$
    \item \tagprob{Scomp Trinomio}{3} $x^4 - 5x^2 + 4$ \textit{(Biquadratica)}
    \item \tagprob{Scomp Trinomio}{3} $x^4 - 13x^2 + 36$
    
    % Ruffini e Misti (10 varianti)
    \item \tagprob{Scomp Ruffini}{2} $x^3 - 2x^2 - x + 2$
    \item \tagprob{Scomp Ruffini}{2} $x^3 - 7x + 6$
    \item \tagprob{Scomp Ruffini}{2} $x^3 - 4x^2 + 5x - 2$
    \item \tagprob{Scomp Ruffini}{2} $x^3 + 2x^2 - 5x - 6$
    \item \tagprob{Scomp Ruffini}{3} $2x^3 - 3x^2 - 8x - 3$
    \item \tagprob{Scomp Ruffini}{3} $3x^3 - 10x^2 + 9x - 2$
    \item \tagprob{Scomp Mista}{3} $x^6 - 64$
    \item \tagprob{Scomp Mista}{3} $x^5 - x^4 - x + 1$
    \item \tagprob{Scomp Mista}{3} $a^4 + 4$ \textit{(Suggerimento: completa il quadrato)}
    \item \tagprob{Scomp Mista}{3} $x^4 + x^2 + 1$
\end{enumerate}
\end{multicols}

\subsection{Frazioni Algebriche: C.E., Semplificazione e Operazioni}
\begin{multicols}{2}
\begin{enumerate}
    % C.E. e Semplificazioni (10 varianti)
    \item \tagprob{Frazioni CE}{1} Determina le C.E.: $\frac{5x}{x^2 - 25}$
    \item \tagprob{Frazioni CE}{1} Determina le C.E.: $\frac{x-1}{x^3 - x}$
    \item \tagprob{Frazioni CE}{2} Determina le C.E.: $\frac{1}{x^2 - 5x + 6}$
    \item \tagprob{Frazioni CE}{2} Determina le C.E.: $\frac{x+y}{x^3 + y^3}$
    \item \tagprob{Fraz Semplifica}{1} Semplifica: $\frac{6a^2b^3}{8a^3b}$
    \item \tagprob{Fraz Semplifica}{1} Semplifica: $\frac{x^2 - 4x + 4}{x^2 - 4}$
    \item \tagprob{Fraz Semplifica}{1} Semplifica: $\frac{a^2 - ab}{a^2 - b^2}$
    \item \tagprob{Fraz Semplifica}{2} Semplifica: $\frac{x^3 - 8}{x^2 - 4}$
    \item \tagprob{Fraz Semplifica}{2} Semplifica: $\frac{2x^2 - 5x + 3}{x^2 - 1}$
    \item \tagprob{Fraz Semplifica}{3} Semplifica: $\frac{x^4 - y^4}{x^3 + x^2y + xy^2 + y^3}$
    
    % Operazioni (10 varianti)
    \item \tagprob{Fraz Addizione}{1} Calcola: $\frac{2}{x-1} + \frac{3}{x+1}$
    \item \tagprob{Fraz Addizione}{1} Calcola: $\frac{x}{x^2 - 4} - \frac{1}{x-2}$
    \item \tagprob{Fraz Addizione}{2} Calcola: $\frac{a+1}{a-2} + \frac{a-1}{a+2} - \frac{4a}{a^2-4}$
    \item \tagprob{Fraz Addizione}{2} Calcola: $\frac{1}{x^2-3x+2} + \frac{1}{x^2-5x+6}$
    \item \tagprob{Fraz Moltiplica}{1} Calcola: $\frac{x^2 - 9}{x^2} \cdot \frac{x}{x - 3}$
    \item \tagprob{Fraz Moltiplica}{2} Calcola: $\frac{a^2 - 4a + 4}{a^2 - a - 2} \cdot \frac{a^2 - 1}{a^2 - 2a}$
    \item \tagprob{Fraz Divisione}{1} Calcola: $\frac{x^2 - y^2}{2x} \div \frac{x + y}{4x^2}$
    \item \tagprob{Fraz Divisione}{2} Calcola: $\frac{x^3 - 1}{x^2 + x} \div \frac{x^2 + x + 1}{x^2 - 1}$
    \item \tagprob{Fraz Mista}{3} Calcola: $\left( \frac{1}{x-1} - \frac{1}{x+1} \right) \cdot \frac{x^2 - 1}{2}$
    \item \tagprob{Fraz Mista}{3} Calcola: $\left( \frac{x+y}{x-y} - \frac{x-y}{x+y} \right) \div \frac{4xy}{x^2 - 2xy + y^2}$
\end{enumerate}
\end{multicols}